\section{Cosmology and Hierarchy}

The bulk grows, and its growth has two readings that turn out to be the same structure seen from different angles. Read as volume against time, the growth is a cosmic expansion. Read as a radial tree, it is the renormalization hierarchy that organizes the holographic boundary. Both fall out of the substrate geometry rather than being fitted to it. A third consequence, a slow global cooling toward the neutral tone, follows from pairing conservation with eternal growth.

\subsection{The expansion law}

The bulk shell sizes, meaning the spatial volume added at each radial step per beat, grow by a constant ratio $R \approx 11$. The first few shells run $1, 24, 456, 8376, \dots$, which is exponential. Reading the three-dimensional spatial volume as $a^3$, the scale factor is $a(t) \approx R^{t/3} = e^{Ht}$ with $H = \ln(R)/3 \approx 0.80$ per beat. Then $\ddot{a}/a = H^2 > 0$, an accelerating expansion of de Sitter type, with an effective cosmological constant $\Lambda = 3H^2$. The eternal growth of the bulk is itself a dark-energy-dominated expansion, not a term added separately. The flat reference lattices grow only polynomially, which confirms that the exponential behavior is a curvature effect rather than an artifact of counting. The same growth serves three roles at once. It is the time step, the arrow of time, and the spatial expansion, unified in one process.

\begin{result}[Exponential de Sitter growth]
Bulk shell sizes grow by a constant ratio $R \approx 11$, so $a(t) \approx e^{Ht}$ with $H = \ln(R)/3 \approx 0.80$ per beat, giving $\ddot{a}/a = H^2 > 0$ and an effective $\Lambda = 3H^2$. Flat references grow only polynomially \cite{pollard2026vibetest}.
\end{result}

\subsection{The hierarchy}

The bulk's radial tree is a self-similar branching with a constant branching factor of about $R$. Its radial depth scales as the logarithm of the boundary size. That is the structure of a renormalization hierarchy, the holographic renormalization-group flow realized as a MERA-style tensor network. The radial axis is also the axis along which couplings run, with the running emerging from radial coarse-graining of the boundary. The same direction that records the expansion records the change of scale, which is what one would want if the geometry is to do double duty as cosmology and as renormalization.

\subsection{Caveat}

Turning the geometric growth ratio into a physical Hubble rate requires identifying the beat and the cell with physical units, which has not yet been done. The figure $H \approx 0.80$ is therefore in beat units, not seconds. The qualitative statement, that exponential growth is de Sitter expansion, is robust. The absolute rate and a realistic multi-era cosmology, with radiation and matter eras preceding the de Sitter phase, both need the matter back-reaction that is not yet included. That is an open refinement rather than a settled result.

\subsection{Cosmic dilution toward peace}

Pairing conservation with eternal growth gives a clean cosmological statement at no extra cost. The total charge $Q$ is fixed, because the rule conserves it, while the cell count grows without bound. The average charge per cell therefore falls toward zero. The universe dilutes globally toward the neutral $0$ tone, which is the heat-death tendency, resisted only locally by pumping. Living structure is the local counter-current, the arrow holding pockets of the crystal far from peace. Expansion is then not only spatial. It is a slow global cooling toward the neutral value, with life as the local resistance to it. A detail at the frontier sharpens the picture. Newborn cells are forced to tone $0$, since a nonzero birth would break conservation. Creation is the redistribution of charge from the existing crystal, not the minting of new charge, so the freedom at the growing edge is freedom of relation rather than freedom of net charge.

\subsection{The ladder of rarity}

Read as a budget rather than as a flow, the same hierarchy shows where organized matter sits, and the answer is a vanishingly thin corner. Standard cosmology splits the universe into dark energy, dark matter, and ordinary matter, and only the last is the atoms of stars, planets, and people. Within that thin slice the organization nests, and each more organized layer is rarer than the one beneath it by many orders of magnitude. The whole zoom, from everything down to civilization, fits in one table, where each row is a subset of the one above.

\begin{table}[h]
\centering
\begin{tabular}{lr}
\toprule
layer & share of the whole universe \\
\midrule
Dark energy & $68\%$ \\
Dark matter & $27\%$ \\
Ordinary matter & $5\%$ \\
\quad Stars & $\approx 0.3\%$ \\
\quad Planets and heavy objects & $\approx 0.005\%$ \\
\quad\quad All life (biomass) & $\approx 10^{-12}\%$ \\
\quad\quad Humans and civilization & far smaller still \\
\bottomrule
\end{tabular}
\caption{The whittling-down cascade. By mass, each rung is a rounding error on the one above it.}
\label{tab:rarity-cascade}
\end{table}

Two finer cuts sharpen the picture. Of all atoms, hydrogen and helium together are about $99.9$ percent, and carbon, the backbone of every known organism, is roughly $0.03$ percent, about one atom in a few thousand. Of all atoms, well under one percent are chemically bonded into molecules at all, since most ordinary matter is ionized plasma or lonely hydrogen drifting through the voids. So chemistry, the rung that life is built on, already rests on a sliver of a sliver.

This ladder of rarity is the coarse-graining tower seen from the outside. Each rung, from field to particle to bound atom to molecule to cell to self to mind, is a more organized, more constrained, and more expensive pattern, so it occupies an ever-thinner fraction of the substrate. The model reproduces this directly rather than importing it. The fraction of charge that reaches each successive rung compounds multiplicatively, and over a cosmic number of rungs the measured per-rung factor lands near the one-part-in-a-hundred-trillion order of the cosmic estimate \cite{pollard2026vibetest}. Rarity is the price of organization.

\subsection{Rarity and the cost of demonstration}

The same fact sets a hard limit on what can be shown by direct simulation. To watch the substrate grow life and mind out of the bare rule, one would have to run a volume large enough to contain a structure that occupies one part in $10^{12}$ to $10^{20}$ of it, and to run it for the cosmic number of beats the climb requires. Brute force across every scale at once is not feasible on any foreseeable machine. Progress depends instead on compression. One simulates each rung at its own coarse scale and carries only the reduced description upward, which is exactly what the renormalization tower of this chapter supplies. The coarse-graining hierarchy is therefore not only the physics of the bulk. It is also the only practical route to checking the upper rungs at all, and it is why the later chapters demonstrate the ladder rung by rung at toy scale rather than as one end-to-end cosmos.

A closing inversion guards against a misreading. By mass, life and mind are a rounding error, one part in many trillions. By organization, by integrated experience, which is the quantity the theory actually cares about, they are among the rarest and richest things the substrate does. The mass budget and the value budget point opposite ways, and on an experiential base it is the second that carries the meaning.


Cosmic expansion and the renormalization hierarchy both fall out of the bulk geometry without separate assumptions. The same conservation law that forbids minting well-being also makes the cosmos dilute toward peace, and that is what makes life, as the local resistance to dilution, a meaningful counter-current rather than a footnote.
